%! @@TITLE@@ — Unit @@UNITINT@@, Lesson @@LESSONID@@ — Experience (Activity · QuickNotes · Application · Check Your Understanding)
% EFFL component, Math Medic sizing: 12pt body, OPEN white space after every question (no
% write-lines). \answerspace{H}{} reserves exactly H of space in the blank; the key passes the
% answer as the second argument, so the two files paginate identically by construction.
% Sizing: 1.4cm ~ 2 handwritten lines, 2.0cm ~ 3, 2.6cm ~ 4. Keep this macro byte-identical
% in the blank and the key. \boxguard counts here are 12pt-relative: use ~14-16.
\documentclass[12pt]{article}
\usepackage{@@PREFIX@@-article}
\usepackage{@@PREFIX@@-boxes}
\newcommand{\answerspace}[2]{\par\nopagebreak\noindent\begin{minipage}[t][#1][t]{\linewidth}%
  \color{keyred}\bfseries #2\end{minipage}\par}

\begin{document}

\pageheader{Unit @@UNITINT@@, Lesson @@LESSONID@@}{Experience: TODO Activity Title}
% NAMESTRIP: no \namedateperiod here — the name row lives on the cover only.

\vspace{0.06in}

% ── Part 1: the Activity ─────────────────────────────────────────────────────
% TWO scenarios, ~10-13 lettered sub-questions total, ~2 pages — the 20-minute timebox.
% Students work from PRIOR KNOWLEDGE ONLY: never name the formal vocabulary here (the
% debrief attaches it). All graphs pre-drawn; students label, circle, and answer in the
% open space. Scenario 2 carries the lesson's crux question (the target misconception).
\begin{headlinebox}{forestbg}
\textbf{TODO Activity Title.} TODO: one motivating context, 2--3 sentences, ending with what
groups should be ready to explain (``how you know \dots\ from the graph'').
\end{headlinebox}

\vspace{0.04in}

\begin{scenariobox}[1. TODO Scenario One]{forest}
TODO: scenario setup. Builds the whole toolkit on one example.
\begin{enumerate}[label=\textbf{\alph*.}, leftmargin=1.8em, itemsep=8pt]
  \item TODO: a table to complete / a rule to write ($f(x) = $ \blank{3.4cm}).
  \item TODO: a feature to circle on the pre-drawn graph, then interpret:
        \answerspace{2.0cm}{}
\end{enumerate}
\end{scenariobox}

\boxguard[16]
\begin{scenariobox}[2. TODO Scenario Two]{forest}
TODO: the contrast case. Carries the crux question, then closes with story-vs-math.
\begin{enumerate}[label=\textbf{\alph*.}, leftmargin=1.8em, itemsep=8pt]
  \item TODO: the crux question (surfaces the target misconception):
        \answerspace{2.6cm}{}
\end{enumerate}
\end{scenariobox}

% ── Part 2: QuickNotes (the debrief fills this) ──────────────────────────────
\boxguard[16]
\begin{tcolorbox}[enhanced, colback=sky, colframe=navy, boxrule=0.8pt, arc=2mm,
    left=3mm, right=3mm, top=2mm, bottom=2mm, breakable,
    fonttitle=\bfseries\small\color{white}, title={QuickNotes: TODO Topic},
    attach boxed title to top left={yshift=-2mm, xshift=4mm},
    boxed title style={colback=navy, colframe=navy, arc=1.5mm}]
% A worked-example figure beside fill-in bullets for the formal terms. A summary of what
% the groups discovered, not a lecture — one box, ~1 page.
\begin{itemize}[leftmargin=1.1em, itemsep=8pt]
  \item TODO: formal term --- definition with \blank{2.4cm} fills.
\end{itemize}
\end{tcolorbox}

% ── Part 3: Practice — Check Your Understanding ──────────────────────────────
\boxguard[16]
\begin{notesbox}{Practice: Check Your Understanding}
Work with your partner, then finish on your own.
\begin{enumerate}[leftmargin=1.7em, itemsep=10pt]
  \item TODO: 4--5 items in new contexts. Include the lesson's one \texttt{work} block and
        close with an SOL-style multiple-choice item (the formative check).
        \answerspace{1.6cm}{}
\end{enumerate}
\end{notesbox}

\end{document}
